El crédito bancario total al sector privado presenta una relación lineal positiva con el Producto Interno Bruto (PIB), evidenciando que conforme la economía ha crecido durante el período analizado, también lo ha hecho el financiamiento otorgado al sector privado. Este comportamiento refleja la estrecha relación existente entre la actividad económica y la expansión del crédito, ya que ambos indicadores tienden a evolucionar de manera conjunta en escenarios de crecimiento económico.
No obstante, se observa un punto atípico correspondiente al segundo trimestre de 2020, período en el que la pandemia de COVID-19 provocó una fuerte contracción de la actividad económica. Las restricciones sanitarias, la disminución de la producción y los cambios en las políticas empresariales y financieras alteraron temporalmente la relación histórica entre ambas variables, generando una desviación respecto de la tendencia observada en el resto de la serie.
En términos generales, la relación observada sugiere que el crecimiento del crédito y el desempeño del PIB se encuentran estrechamente vinculados. Por ello, las decisiones de política monetaria y financiera deben considerar el comportamiento conjunto de ambos indicadores, ya que cambios significativos en uno de ellos pueden influir en la dinámica del otro y, en consecuencia, en el desempeño de la economía nacional.